\documentclass[conference]{IEEEtran}
\usepackage[utf8]{inputenc}
\usepackage[portuguese]{babel}
\usepackage{amsmath}
\usepackage{algorithm}
\usepackage{algorithmic}
\usepackage{graphicx}
\usepackage{cite}
\usepackage{url}
\usepackage{enumitem}

\setlist[itemize]{itemsep=3pt, parsep=3pt, topsep=3pt}

\begin{document}

\title{Análise Comparativa de Algoritmos de Ordenação em Estruturas de Dados Dinâmicas e Estáticas}

\author{
\IEEEauthorblockN{Elder Rayan Oliveira Silva, Samuel Wagner Tiburi Silveira, Diogo Gomes Figueiredo, \\
Manoel Junio Duarte da Silva, Pedro Yan Alcantara Palácio, Antonio Pereira da Luz Neto}
\IEEEauthorblockA{
Universidade Federal do Cariri - UFCA\\
Engenharia de Software\\
Disciplina: Estrutura de Dados\\
}
}

\maketitle

\section{Introdução}

Em produção.

% ============================================================
\section{Desenvolvimento}
% ============================================================

Esta seção descreve a implementação prática do sistema de ranking de alunos, detalhando as estruturas de dados utilizadas e os algoritmos de ordenação implementados.

\subsection{Estruturas de Dados}

Para avaliar o impacto da organização da memória no desempenho dos algoritmos, foram implementadas duas estruturas de dados distintas para armazenar o ranking de alunos, ambas utilizando o conceito de \textbf{lista duplamente encadeada}:

\subsubsection{Lista Encadeada Dinâmica}

A lista dinâmica utiliza alocação de memória sob demanda, através das funções \texttt{malloc()} e \texttt{free()} da biblioteca padrão C. Cada nó da lista contém:

\begin{itemize}
    \item \textbf{ID}: Identificador único do aluno
    \item \textbf{Nota}: Nota numérica do aluno (0-100)
    \item \textbf{Ponteiros}: Referências para o próximo e para o nó anterior (\texttt{prox} e \texttt{ant})
\end{itemize}

\textbf{Vantagens:}
\begin{itemize}
    \item Tamanho dinâmico, ajustável conforme necessário
    \item Inserções e remoções eficientes em posições arbitrárias
    \item Navegação bidirecional facilitada pela estrutura duplamente encadeada
    \item Não desperdiça memória com espaços pré-alocados não utilizados
\end{itemize}

\textbf{Desvantagens:}
\begin{itemize}
    \item Acesso sequencial obrigatório (sem acesso direto por índice)
    \item Overhead de memória devido ao dobro de ponteiros por nó
    \item Possível fragmentação de memória
    \item Localidade espacial de cache prejudicada
\end{itemize}

\subsubsection{Lista Encadeada Estática}

A lista estática simula uma lista duplamente encadeada utilizando um vetor pré-alocado de tamanho fixo. Os "ponteiros" são implementados como índices inteiros do vetor (\texttt{prox} e \texttt{ant}), e uma lista de posições livres é mantida para gerenciar a reutilização de espaços.

\textbf{Vantagens:}
\begin{itemize}
    \item Melhor localidade de cache (dados contíguos na memória)
    \item Sem overhead de chamadas a \texttt{malloc/free}
    \item Tempo de alocação previsível e constante
    \item Acesso mais rápido devido à contiguidade da memória
\end{itemize}

\textbf{Desvantagens:}
\begin{itemize}
    \item Tamanho máximo fixo definido em tempo de compilação
    \item Possível desperdício de memória se o vetor não for totalmente utilizado
    \item Complexidade adicional no gerenciamento de índices livres
\end{itemize}

\subsubsection{Justificativa da Escolha}

A implementação de ambas as estruturas permite avaliar empiricamente o trade-off entre flexibilidade e desempenho. A lista dinâmica oferece flexibilidade máxima, ideal para cenários onde o tamanho do dataset é imprevisível. Já a lista estática sacrifica flexibilidade em favor de desempenho, sendo ideal para aplicações críticas onde o limite superior de dados é conhecido.

No contexto de aplicações reais (como sistemas de ranking), a escolha depende de fatores como:
\begin{itemize}
    \item Volume máximo esperado de dados
    \item Frequência de inserções/remoções
    \item Requisitos de latência
    \item Restrições de memória do sistema
\end{itemize}

\subsection{Algoritmos de Ordenação Implementados}

Cinco algoritmos clássicos de ordenação foram implementados e testados em ambas as estruturas de dados:

\subsubsection{Bubble Sort}

Algoritmo baseado em comparações sucessivas de elementos adjacentes, realizando trocas quando estão fora de ordem. Percorre a lista repetidamente até que nenhuma troca seja necessária.

\textbf{Características:}
\begin{itemize}
    \item Algoritmo estável (mantém ordem relativa de elementos iguais)
    \item In-place (não requer memória auxiliar significativa)
    \item Simples de implementar
    \item Ineficiente para grandes volumes de dados
\end{itemize}

\subsubsection{Selection Sort}

A cada iteração, seleciona o menor elemento não ordenado e o posiciona em sua posição final. Divide conceitualmente a lista em duas partes: ordenada e não ordenada.

\textbf{Características:}
\begin{itemize}
    \item Não estável na implementação padrão
    \item In-place
    \item Número fixo de trocas: $O(n)$
    \item Útil quando o custo de troca é alto
\end{itemize}

\subsubsection{Insertion Sort}

Constrói a lista ordenada incrementalmente, inserindo cada elemento na posição correta dentro da porção já ordenada.

\textbf{Características:}
\begin{itemize}
    \item Algoritmo estável
    \item In-place
    \item Adaptativo (eficiente para dados quase ordenados)
    \item Excelente para listas pequenas ou parcialmente ordenadas
\end{itemize}

\subsubsection{Quick Sort}

Algoritmo de divisão e conquista que escolhe um pivô e particiona a lista em elementos menores e maiores que o pivô, ordenando recursivamente cada partição.

\textbf{Características:}
\begin{itemize}
    \item Não estável na implementação padrão
    \item In-place (considerando a pilha de recursão)
    \item Geralmente o mais rápido na prática
    \item Desempenho degradado no pior caso
\end{itemize}

\subsubsection{Merge Sort}

Algoritmo de divisão e conquista que divide recursivamente a lista em sublistas menores, ordenando-as e mesclando os resultados.

\textbf{Características:}
\begin{itemize}
    \item Algoritmo estável
    \item Requer memória auxiliar $O(n)$
    \item Desempenho consistente independente dos dados
    \item Ideal para dados em armazenamento externo
\end{itemize}

\subsection{Algoritmo Ideal: Merge Sort}

Para o contexto deste projeto, o \textbf{Merge Sort} foi selecionado como o algoritmo ideal, baseado nos seguintes critérios:

\subsubsection{Garantia de Desempenho}

O Merge Sort apresenta complexidade $O(n \log n)$ em \textbf{todos os casos} (melhor, médio e pior). Diferentemente do Quick Sort, que pode degradar para $O(n^2)$ no pior caso, o Merge Sort oferece previsibilidade de desempenho crucial para sistemas críticos.

\subsubsection{Estabilidade}

A propriedade de estabilidade é essencial em sistemas de ranking onde múltiplos critérios de ordenação podem ser aplicados sequencialmente. Ao ordenar primeiro por nota e depois por nome, o Merge Sort preserva a ordem relativa de alunos com mesma nota.

\subsubsection{Funcionamento Detalhado}

O algoritmo opera em três fases:

\textbf{1. Divisão (Divide):}
\begin{algorithmic}
\IF{$n \leq 1$}
    \RETURN
\ENDIF
\STATE $meio \gets n / 2$
\STATE Divide lista em $L$ (primeira metade) e $R$ (segunda metade)
\end{algorithmic}

\textbf{2. Conquista (Conquer):}
\begin{algorithmic}
\STATE MergeSort($L$)
\STATE MergeSort($R$)
\end{algorithmic}

\textbf{3. Combinação (Merge):}
\begin{algorithmic}
\STATE Inicializa lista auxiliar $A$
\WHILE{$L$ não vazia \AND $R$ não vazia}
    \IF{$L[0] \leq R[0]$}
        \STATE Adiciona $L[0]$ em $A$ e remove de $L$
    \ELSE
        \STATE Adiciona $R[0]$ em $A$ e remove de $R$
    \ENDIF
\ENDWHILE
\STATE Adiciona elementos restantes de $L$ ou $R$ em $A$
\end{algorithmic}

\textbf{Custo Espacial:}

O principal trade-off do Merge Sort é o uso de memória auxiliar $O(n)$ para a fusão das sublistas. No entanto, este custo é aceitável considerando:
\begin{itemize}
    \item Memória moderna abundante em sistemas computacionais
    \item Possibilidade de otimização para ordenação externa
    \item Benefício da consistência de desempenho
\end{itemize}

\textbf{Adaptações para Listas Encadeadas:}

Para listas encadeadas, o Merge Sort é particularmente eficiente pois:
\begin{itemize}
    \item A divisão pode ser feita alterando ponteiros (sem cópia de dados)
    \item A fusão se torna naturalmente in-place através de religação de nós
    \item Não sofre com o acesso sequencial, pois o acesso direto não é necessário
\end{itemize}

% ============================================================
\section{Análise Teórica}
% ============================================================

Esta seção apresenta a análise assintótica de complexidade computacional dos cinco algoritmos implementados, considerando os cenários de melhor, médio e pior caso.

\subsection{Notação Big-O e Análise de Complexidade}

A notação Big-O descreve o comportamento assintótico de um algoritmo à medida que o tamanho da entrada $n$ cresce. As análises consideram:
\begin{itemize}
    \item \textbf{Complexidade de Tempo:} Número de operações elementares
    \item \textbf{Complexidade de Espaço:} Memória auxiliar necessária
\end{itemize}

\subsection{Bubble Sort}

\subsubsection{Melhor Caso: $O(n)$}

Ocorre quando a lista já está ordenada. Com a otimização de flag de troca, o algoritmo detecta a ordenação após uma única passagem.

\textbf{Análise:}
\begin{equation}
T_{melhor}(n) = n - 1 \text{ comparações} = O(n)
\end{equation}

\subsubsection{Caso Médio: $O(n^2)$}

Para dados aleatórios, em média metade das comparações resulta em trocas.

\textbf{Análise:}
\begin{equation}
T_{medio}(n) = \sum_{i=1}^{n-1} i = \frac{n(n-1)}{2} = O(n^2)
\end{equation}

\subsubsection{Pior Caso: $O(n^2)$}

Ocorre quando a lista está inversamente ordenada. Toda comparação resulta em troca.

\textbf{Análise:}
\begin{equation}
T_{pior}(n) = \sum_{i=1}^{n-1} i = \frac{n(n-1)}{2} = O(n^2)
\end{equation}

\textbf{Complexidade Espacial:} $O(1)$ - in-place

\subsection{Selection Sort}

\subsubsection{Todos os Casos: $O(n^2)$}

O Selection Sort executa o mesmo número de comparações independentemente da configuração inicial dos dados.

\textbf{Análise:}
\begin{equation}
T(n) = \sum_{i=1}^{n-1} (n-i) = \frac{n(n-1)}{2} = O(n^2)
\end{equation}

\textbf{Justificativa:}
Para cada posição $i$, o algoritmo deve percorrer os $(n-i)$ elementos restantes para encontrar o mínimo, resultando sempre em $\frac{n(n-1)}{2}$ comparações.

\textbf{Número de Trocas:}
Diferentemente de Bubble Sort, o número de trocas é sempre $O(n)$ (no máximo uma por posição), tornando-o vantajoso quando o custo de troca é significativamente maior que o de comparação.

\textbf{Complexidade Espacial:} $O(1)$ - in-place

\subsection{Insertion Sort}

\subsubsection{Melhor Caso: $O(n)$}

Ocorre quando a lista já está ordenada. Cada elemento requer apenas uma comparação.

\textbf{Análise:}
\begin{equation}
T_{melhor}(n) = n - 1 \text{ comparações} = O(n)
\end{equation}

\subsubsection{Caso Médio: $O(n^2)$}

Para dados aleatórios, cada elemento é inserido, em média, no meio da porção ordenada.

\textbf{Análise:}
\begin{equation}
T_{medio}(n) = \sum_{i=1}^{n-1} \frac{i}{2} = \frac{n(n-1)}{4} = O(n^2)
\end{equation}

\subsubsection{Pior Caso: $O(n^2)$}

Ocorre quando a lista está inversamente ordenada. Cada elemento deve ser comparado com todos os anteriores.

\textbf{Análise:}
\begin{equation}
T_{pior}(n) = \sum_{i=1}^{n-1} i = \frac{n(n-1)}{2} = O(n^2)
\end{equation}

\textbf{Propriedade Adaptativa:}

O Insertion Sort é adaptativo, ou seja, seu desempenho melhora significativamente para dados parcialmente ordenados. Para uma lista com $k$ inversões:
\begin{equation}
T(n) = O(n + k)
\end{equation}

\textbf{Complexidade Espacial:} $O(1)$ - in-place

\subsection{Quick Sort}

\subsubsection{Melhor Caso: $O(n \log n)$}

Ocorre quando o pivô divide consistentemente a lista em duas partes aproximadamente iguais.

\textbf{Análise:}
Profundidade da recursão: $\log_2 n$\\
Trabalho por nível: $O(n)$ (particionamento)
\begin{equation}
T_{melhor}(n) = O(n) \times \log_2 n = O(n \log n)
\end{equation}

\subsubsection{Caso Médio: $O(n \log n)$}

Para dados aleatórios, o pivô divide a lista em proporções variadas, mas a profundidade esperada permanece $O(\log n)$.

\textbf{Análise (Relação de Recorrência):}
\begin{equation}
T(n) = T(k) + T(n-k-1) + O(n)
\end{equation}
onde $k$ é o número de elementos menores que o pivô.

Para pivô mediano (caso esperado):
\begin{equation}
T_{medio}(n) = 2T(n/2) + O(n) = O(n \log n)
\end{equation}

\subsubsection{Pior Caso: $O(n^2)$}

Ocorre quando o pivô é sempre o menor ou maior elemento (lista já ordenada ou inversamente ordenada com pivô ingênuo).

\textbf{Análise:}
Profundidade da recursão: $n$\\
Partições extremamente desbalanceadas
\begin{equation}
T_{pior}(n) = T(n-1) + O(n) = \sum_{i=1}^{n} i = O(n^2)
\end{equation}

\textbf{Estratégias de Mitigação:}
\begin{itemize}
    \item Escolha aleatória de pivô: reduz probabilidade do pior caso
    \item Mediana-de-três: pivô = mediana(primeiro, meio, último)
    \item Hybrid approach: usar Insertion Sort para sublistas pequenas
\end{itemize}

\textbf{Complexidade Espacial:} $O(\log n)$ (pilha de recursão no caso médio), $O(n)$ no pior caso

\subsection{Merge Sort}

\subsubsection{Todos os Casos: $O(n \log n)$}

O Merge Sort apresenta complexidade consistente independente da configuração dos dados.

\textbf{Análise (Relação de Recorrência):}
\begin{equation}
T(n) = 2T(n/2) + O(n)
\end{equation}

Aplicando o Teorema Mestre (Caso 2):
\begin{equation}
T(n) = O(n \log n)
\end{equation}

\textbf{Demonstração Detalhada:}

Profundidade da árvore de recursão: $\log_2 n$\\
Trabalho por nível: $O(n)$ (fusão de todas as sublistas)

Nível 0: $n$ elementos em 1 lista\\
Nível 1: $n$ elementos em 2 listas\\
Nível 2: $n$ elementos em 4 listas\\
...\\
Nível $\log n$: $n$ elementos em $n$ listas

\begin{equation}
T(n) = \sum_{i=0}^{\log n} n = n \times (\log n + 1) = O(n \log n)
\end{equation}

\textbf{Invariância do Desempenho:}

O Merge Sort sempre divide a lista ao meio, independente dos valores. Portanto:
\begin{align}
T_{melhor}(n) &= O(n \log n) \\
T_{medio}(n) &= O(n \log n) \\
T_{pior}(n) &= O(n \log n)
\end{align}

\textbf{Complexidade Espacial:} $O(n)$ - requer array auxiliar para fusão

\subsection{Comparação Teórica}

A Tabela \ref{tab:complexidade} resume as complexidades dos cinco algoritmos:

\begin{table}[h]
\centering
\caption{Complexidade de Tempo dos Algoritmos de Ordenação}
\label{tab:complexidade}
\renewcommand{\arraystretch}{1.3}
\begin{tabular}{|l|c|c|c|}
\hline
\textbf{Algoritmo} & \textbf{Melhor} & \textbf{Médio} & \textbf{Pior} \\
\hline
Bubble Sort & $O(n)$ & $O(n^2)$ & $O(n^2)$ \\
Selection Sort & $O(n^2)$ & $O(n^2)$ & $O(n^2)$ \\
Insertion Sort & $O(n)$ & $O(n^2)$ & $O(n^2)$ \\
Quick Sort & $O(n \log n)$ & $O(n \log n)$ & $O(n^2)$ \\
Merge Sort & $O(n \log n)$ & $O(n \log n)$ & $O(n \log n)$ \\
\hline
\end{tabular}
\end{table}

\subsection{Considerações Teóricas}

\subsubsection{Algoritmos $O(n^2)$}

Os três primeiros algoritmos (Bubble, Selection, Insertion) compartilham complexidade quadrática no caso médio e pior, mas apresentam diferenças práticas importantes:

\begin{itemize}
    \item \textbf{Insertion Sort} é adaptativo e eficiente para dados parcialmente ordenados
    \item \textbf{Selection Sort} minimiza o número de trocas (útil quando trocas são custosas)
    \item \textbf{Bubble Sort} é raramente usado na prática devido à sua ineficiência geral
\end{itemize}

\subsubsection{Algoritmos $O(n \log n)$}

Os algoritmos eficientes apresentam trade-offs distintos:

\begin{itemize}
    \item \textbf{Merge Sort:} Consistente e estável, mas requer $O(n)$ de espaço
    \item \textbf{Quick Sort:} Mais rápido na prática (boas constantes), mas instável e pode degradar para $O(n^2)$
\end{itemize}

\subsubsection{Limite Inferior Teórico}

Qualquer algoritmo de ordenação baseado em comparações requer pelo menos:
\begin{equation}
\Omega(n \log n) \text{ comparações no pior caso}
\end{equation}

\textbf{Prova (Informação Teórica):}

Existem $n!$ permutações possíveis de $n$ elementos. Uma árvore de decisão binária requer profundidade mínima de:
\begin{equation}
h \geq \log_2(n!) = \Omega(n \log n)
\end{equation}

Portanto, Merge Sort e Quick Sort (caso médio) são assintoticamente ótimos.

\subsection{Implicações Práticas}

\subsubsection{Escolha do Algoritmo}

A seleção do algoritmo ideal depende de múltiplos fatores:

\begin{enumerate}
    \item \textbf{Tamanho do Dataset:}
    \begin{itemize}
        \item $n < 50$: Insertion Sort (simplicidade e overhead baixo)
        \item $n$ médio/grande: Quick Sort ou Merge Sort
    \end{itemize}

    \item \textbf{Requisitos de Estabilidade:}
    \begin{itemize}
        \item Necessária: Merge Sort ou Insertion Sort
        \item Dispensável: Quick Sort (melhor desempenho prático)
    \end{itemize}

    \item \textbf{Restrições de Memória:}
    \begin{itemize}
        \item Crítica: Quick Sort ou algoritmos in-place
        \item Flexível: Merge Sort
    \end{itemize}

    \item \textbf{Características dos Dados:}
    \begin{itemize}
        \item Parcialmente ordenados: Insertion Sort
        \item Aleatórios: Quick Sort
        \item Distribuição desconhecida: Merge Sort (garantia de desempenho)
    \end{itemize}
\end{enumerate}

\subsubsection{Otimizações Híbridas}

Na prática, bibliotecas de ordenação modernas utilizam abordagens híbridas:

\begin{itemize}
    \item \textbf{Introsort (C++ std::sort):} Quick Sort com fallback para Heap Sort quando detecta $O(n^2)$
    \item \textbf{Timsort (Python, Java):} Merge Sort adaptativo que explora sequências ordenadas
    \item \textbf{Hybrid Quick/Insertion:} Quick Sort para partições grandes, Insertion para pequenas
\end{itemize}

% ============================================================
% Conclusão desta seção (se necessária)
% ============================================================

As análises teóricas apresentadas fornecem a base para a compreensão dos resultados experimentais que serão apresentados nas seções subsequentes. A confirmação empírica das complexidades teóricas, bem como a avaliação do impacto das estruturas de dados, são objetivos centrais deste trabalho.

\end{document}
